% ---- BEGIN monograph/chapters/14_discussion.tex ----
\chapter{Discussion}

\noindent\textbf{Status:} \texttt{INTEGRATED THEORY STATE / OPEN FRONTIERS EXPLICIT / MONOGRAPH TEMPORAL SPINE CLOSED TO EINSTEIN INTERFACE}.

The construction developed in this monograph treats temporal structure as an ordered relational process conditionally constructed from the declared IDT primitives, with its own elapsed measure, phase transport, present-event support, bifurcation structure, transport law, memory substrate, and reconstructive geometry. It does not claim that physical time follows uniquely from information alone. The final chapters then place that temporal branch beside an independently typed spatial/source branch and carry both to an Einstein--Bianchi closure test.

The resulting dependency chain is
\[
\boxed{
\begin{gathered}
\mathrm{TIR}
\rightarrow
\mathrm{Temporal\ Primitive}
\rightarrow
\mathrm{Temporal\ Wave}
\rightarrow
\mathrm{NOW}
\rightarrow
\mathrm{Bifurcation}
\\
{}\rightarrow
\mathrm{Temporal\ Transport}
\rightarrow
\mathrm{Memory}
\rightarrow
\mathrm{ORCHORBITAL}
\\
{}\rightarrow
\mathrm{Retrodiction}
\rightarrow
\mathrm{Retrocausal\ Tests}
\rightarrow
\mathrm{Einstein\ Closure}
\end{gathered}
}.
\]
A parallel source chain runs from the temporal primitive through gauge-covariant Noether/source interfaces and the relativistic field bridge to the same Einstein closure.

\section{Time as positive relational activity}

A central simplification of the theory is the decomposition of directed transition rates into positive activity and signed current,
\[
\mathfrak a=W_++W_->0,
\qquad
\mathfrak j=W_+-W_-.
\]
The positive part generates elapsed measure,
\[
d\Theta=\mathfrak a\,d\lambda,
\]
while the ratio of local and reference activities generates the relational lapse,
\[
N_R(x|r)=\frac{\mathfrak a_x}{\mathfrak a_r}.
\]
This produces an ordering-compatible clock architecture before a spacetime metric is introduced. Reparameterisation invariance follows from the one-density transformation of activity, and reference changes compose multiplicatively.

The signed component retains directional information. The edge affinity is recovered through
\[
A=2\operatorname{artanh}\left(\frac{\mathfrak j}{\mathfrak a}\right)
\]
for finite positive rates. Thus pace and direction are related but typed separately. This distinction propagates through the later theory: positive support defines admitted NOW events, while oriented phase and winding variables retain directional lineage.

\section{Relative information and clock geometry}

The local memoryless clock realisation converts the lapse into a Shannon relative-information scalar,
\[
\Phi(N_R)=N_R-1-\ln N_R.
\]
Its convexity and Fisher Hessian show that the neighbourhood of a reference clock already carries information geometry. This gives a concrete bridge between the temporal primitive and the broader relational information geometry used elsewhere in TIR/RFC.

The physical significance is experimentally addressable. A clock ratio may be measured directly after calibration, while \(\Phi_N=c^2\ln N_R\) supplies a weak-field comparison coordinate. The strongest version of the bridge is therefore an equality between two independently generated quantities: the temporal activity ratio and the clock ratio implied by the admitted gravitational geometry.

\section{Memory changes the temporal question}

Once event imprints are persisted, the theory can ask a stronger question than whether a system evolves forward. It can ask which information about earlier transitions remains reconstructible from a later state plus a finite lineage carrier.

The memory reference class makes this question constructive. Smooth segments admit reversible Kepler propagation, event kicks have explicit inverses, residence cells commit chronology and active-attractor identity, and winding/radius coordinates provide a compressed lift back to the complete position lineage. On the declared positive-radius strata,
\[
(\alpha,\mathcal W,\rho_1,\ldots,\rho_{N-1},r_N)
\longrightarrow
(r_1,\ldots,r_N)
\longrightarrow
(u_1,\ldots,u_N)
\]
is an explicit reconstruction chain.

The scalar budget result
\[
N_{\rm radial}=\frac12N_{\rm Cartesian}
\]
shows that temporal lineage can be reconstructively sufficient without retaining every Cartesian coordinate of the reference trajectory. The residence binding adds provenance to that compression by connecting each radius to the event cell and post-state commitment that generated it.

The current mathematical frontier is global-domain coverage. The stratified decoder provides exact constructive results on the declared nonsingular domain; the remaining programme is to classify singular strata and any residual collisions until the global injectivity question is exhausted.

\section{Retrodiction and causal interpretation}

The distinction between retrodiction and retrocausal evidence is one of the main methodological outcomes of the programme. E003 demonstrates how later checkpoints can identify withheld variables in an explicitly invertible dynamical model. This is a property of information retained by the evolution and observation map.

A future-conditioned physical claim has a different evidential structure. Its decisive observable is a distributional dependence of an already committed earlier record on a condition selected later,
\[
P(X_{\mathrm{past}}\mid Y_{\mathrm{future}})
\stackrel{?}{\neq}
P(X_{\mathrm{past}}),
\]
together with cryptographic chronology, blinding, channel audits, sham conditions, and replication. The theory therefore provides both a mechanism for reconstructive inference and a firewall that prevents reconstructive success from being promoted automatically into a causal claim.

\section{Why space is introduced late}

The temporal branch derives an elapsed measure and clock comparison without using spatial geometry as its primitive definition. The spatial branch remains an independent TIR development until the relativistic interface. At that point the temporal one-form
\[
\Theta_R=N_Rc\,dt
\]
and the spatial triad \(e^{\hat i}\) can be assembled into a spacetime coframe.

This ordering gives the spacetime phase relation
\[
\Delta\phi
=
\frac{1}{\hbar}
(\mathbf p\cdot\Delta\mathbf x-E\Delta t)
\]
a genuine closure role. The relation tests the recombined branches. It therefore becomes a diagnostic of the temporal map, the spatial map, and their interface rather than an upstream assumption from which the temporal branch is simply read off.

\section{Einstein closure as a conservation problem}

At the relativistic endpoint, conservation is the organising criterion. The geometric identity
\[
\nabla^\mu G_{\mu\nu}=0
\]
requires the complete admitted source side to have matching divergence. The Maxwell--matter bridge supplies
\[
\nabla^\mu T^{\rm EM}_{\mu\nu}
=-F_{\nu\lambda}J_{\rm EM}^{\lambda},
\qquad
\nabla^\mu T^{\rm matter}_{\mu\nu}
=+F_{\nu\lambda}J_{\rm EM}^{\lambda},
\]
so that the total stress tensor is conserved on the admitted equations.

This establishes the correct structural interface for
\[
G_{\mu\nu}+\Lambda g_{\mu\nu}
=
\frac{8\pi G}{c^4}T^{\rm total}_{\mu\nu}.
\]
The current research programme now concentrates on universality and source composition: the same gravitational normalisation must survive distinct physical systems, all relevant matter sectors must enter the stress ledger, and any dynamic vacuum term must carry its own conservation-compatible source model.

\section{Independent physical channels}

Several downstream channels provide opportunities to challenge the closure from different directions.

The first is precision clock comparison, where \(N_R\) is the common observable between temporal activity and gravitational geometry. The second is spectroscopy, where discrete energy differences generate temporal relative-phase lines and the half-interface supplies sharply specified interference nulls in compatible two-channel systems. The third is the neutrino directional-stress programme, where response-calibrated angular data can be converted into TT stress observables and compared with the spin-2 source model. The fourth is horizon thermodynamics, where the Euclidean \(2\pi\) winding and spin-half antiperiodicity provide exact boundary conditions.

These channels share a useful property: each exposes a different part of the dependency graph. Agreement across several channels would constrain the interfaces far more strongly than success in one numerical benchmark.

\section{Recovery is not discrimination}

The final architecture distinguishes standard-physics recovery from novel empirical discrimination. Recovery of the quantum phase-frequency relation, Einstein--Bianchi conservation, linearised TT propagation, or Hawking/Gibbons--Hawking horizon thermodynamics is necessary for compatibility, but those established relations are not counted as new IDT predictions. A novel physical claim begins only when an IDT quantity is fixed independently of the target observation and can disagree with a declared baseline on the same held-out observable.

Chapter~13 therefore promotes three explicit discriminator classes: an independently generated activity-to-clock predictor, a physically calibrated half-interface \(4\pi\) test, and the preregistered future-conditioned committed-record protocol. Failure of any one of these is informative because the calibration map, null, and promotion rule are defined before the target result is inspected.

\section{Evidence architecture}

The project uses evidence classes because mathematical and physical statements mature at different rates. A useful hierarchy is
\[
\text{definition}
\prec
\text{proved identity}
\prec
\text{reference-model result}
\prec
\text{calibrated physical correspondence}
\prec
\text{replicated measured result}.
\]
The exact identities in the clock, information-geometric, memory, and horizon sectors can therefore coexist with active physical calibration gates in the gravitational and retrocausal sectors without ambiguity about their status.

The same principle applies to computational evidence. A hosted test count is attached to a commit and test coordinate, but the count is a software-assurance receipt rather than a multiplicity of physical observations. A content hash binds a memory packet to its lineage; an FPDG failure receipt maps a failed test to the precise claim/interface coordinate when such a binding exists. Provenance is consequently part of the theory's operational semantics, especially in a programme whose subject is temporal order itself.

\section{Current frontier}

The monograph closes the conceptual temporal spine at the Einstein interface, while several research coordinates remain active:
\begin{enumerate}
\item completion of global-domain injectivity for the retrodiction quotient/fibre construction;
\item execution and independent replication of the preregistered physical future-conditioned protocol;
\item precision binding of the relational lapse to physical clocks and gravitational redshift;
\item full source composition and cross-system calibration of the Einstein coupling;
\item response-corrected IceTracks directional-stress analysis;
\item a first-principles dynamical source model for any promoted spacetime-dependent vacuum term;
\item unified-limit audits joining the temporal, electromagnetic, matter, and gravitational ledgers.
\end{enumerate}

These are constructive frontiers: each is represented by an explicit observable, collision condition, conservation equation, calibration factor, or preregistered statistical test.

\section{Conclusion}

Informational Dynamics of Time reaches a coherent endpoint by treating temporal order as a structure conditionally derived from explicitly declared relational primitives and carrying that structure through phase, NOW, bifurcation, transport, memory, reconstructive inference, and finally spacetime conservation. Its strongest established results are the exact algebraic and constructive identities inside those declared domains; its strongest physical programme is the set of experiments that compare independently generated temporal observables with clock, phase, stress, and geometric measurements.

The resulting architecture can be summarised by one closure principle:
\[
\boxed{
\text{time is derived first; spacetime is tested when time and space are recombined.}
}
\]
This principle keeps the temporal construction identifiable, makes its failure modes localisable, and turns Einstein closure into an empirical and mathematical endpoint rather than an assumed starting geometry.

% ---- END monograph/chapters/14_discussion.tex ----

% ---- BEGIN monograph/appendices/A_harmonic_admission_scheduler.tex ----
\chapter[Harmonic Admission Scheduling]{Harmonic Admission Scheduling for Computational Experiments}
\label{app:harmonic-admission}

This appendix records an engineering control used while running the computational parts of the research program. Its evidential class is \texttt{ENGINEERING\_CONTROL}: it regulates when local processes may start so that numerical work, rendering, repository I/O, and document compilation are less likely to collide in the same execution window. The scheduling convention does not participate in the derivation of the temporal field equations.

\section{Dyadic admission bands}

A base cadence
\[
f_0 = 7.83\;\mathrm{Hz},
\qquad
T_0=f_0^{-1},
\]
is used as an implementation clock. The workload bands are defined by
\[
f_k=\frac{f_0}{2^k},
\qquad
T_k=\frac{2^k}{f_0}.
\]
Band $B_0$ is reserved for lightweight heartbeat observation. Workload admission begins at $B_1$.

For integer base tick $n$, the phase class of band $B_k$, $k\geq1$, is
\[
n\equiv 2^{k-1}\pmod{2^k}.
\]
These classes are pairwise disjoint. Consequently, a base tick can admit at most one workload class among $B_1,\ldots,B_K$.

\begin{center}
\begin{tabular}{ccl}
\hline
Band & Frequency & Operational role \\
\hline
$B_0$ & $7.830$ Hz & heartbeat observation \\
$B_1$ & $3.915$ Hz & queue and admission \\
$B_2$ & $1.9575$ Hz & candidate microsteps \\
$B_3$ & $0.97875$ Hz & critic and graph audit \\
$B_4$ & $0.489375$ Hz & numerics and small slices \\
$B_5$ & $0.2446875$ Hz & serialization and repository I/O \\
$B_6$ & $0.12234375$ Hz & visualization \\
$B_7$ & $0.061171875$ Hz & LaTeX and heavy builds \\
\hline
\end{tabular}
\end{center}

The numerical value of $f_0$ is an operational scheduling convention in this appendix. Any physical use of a frequency scale requires an independent claim, derivation, and evidence path in the main formalism.

\section{Skip semantics and jitter control}

An admission window is permission to start queued work rather than a demand for periodic execution. Empty windows perform no work. A missed window is skipped:
\[
\text{missed slot}\longrightarrow\text{skip},
\]
so delayed processes cannot produce a catch-up burst. If a task assigned to $B_k$ is still active at its next eligible slot, the new instance is skipped as well.

Measured scheduling jitter $J_k$ can demote a workload to a slower band,
\[
J_k>J_{\max}
\quad\Longrightarrow\quad
B_k\mapsto B_{k+1},
\]
subject to the configured maximum band. This mechanism is intended to protect the interactive research workflow from render, I/O, and compilation spikes while keeping the computational procedures reproducible.

\section{Reference implementation}

The reference implementation is stored in \texttt{src/idt/harmonic\_scheduler.py}. The scheduler exposes the band definitions, exact phase-slot predicate, next-slot calculation, quiet deep-boundary slots, and jitter-triggered demotion. Reference tests verify pairwise disjoint workload classes and the no-catch-up admission semantics.

% ---- END monograph/appendices/A_harmonic_admission_scheduler.tex ----

% ---- BEGIN monograph/frontmatter/references.tex ----
\renewcommand{\bibname}{References}
\addcontentsline{toc}{chapter}{References}
\begin{thebibliography}{99}

\bibitem{Shannon1948}
C.~E. Shannon,
\emph{A Mathematical Theory of Communication},
Bell System Technical Journal, vol.~27, no.~3, pp.~379--423, 1948.
\href{https://doi.org/10.1002/j.1538-7305.1948.tb01338.x}{doi:10.1002/j.1538-7305.1948.tb01338.x}.

\bibitem{KullbackLeibler1951}
S.~Kullback and R.~A. Leibler,
\emph{On Information and Sufficiency},
The Annals of Mathematical Statistics, vol.~22, no.~1, pp.~79--86, 1951.
\href{https://doi.org/10.1214/aoms/1177729694}{doi:10.1214/aoms/1177729694}.

\bibitem{Pancharatnam1956}
S.~Pancharatnam,
\emph{Generalized theory of interference, and its applications. Part I. Coherent pencils},
Proceedings of the Indian Academy of Sciences, Section A, vol.~44, no.~5, pp.~247--262, 1956.
\href{https://doi.org/10.1007/BF03046050}{doi:10.1007/BF03046050}.

\bibitem{Berry1984}
M.~V. Berry,
\emph{Quantal phase factors accompanying adiabatic changes},
Proceedings of the Royal Society of London. A, vol.~392, no.~1802, pp.~45--57, 1984.
\href{https://doi.org/10.1098/rspa.1984.0023}{doi:10.1098/rspa.1984.0023}.

\bibitem{ProvostVallee1980}
J.~P. Provost and G.~Vallee,
\emph{Riemannian structure on manifolds of quantum states},
Communications in Mathematical Physics, vol.~76, no.~3, pp.~289--301, 1980.
\href{https://doi.org/10.1007/BF02193559}{doi:10.1007/BF02193559}.

\bibitem{Onsager1931}
L.~Onsager,
\emph{Reciprocal Relations in Irreversible Processes. I},
Physical Review, vol.~37, pp.~405--426, 1931.
\href{https://doi.org/10.1103/PhysRev.37.405}{doi:10.1103/PhysRev.37.405}.

\bibitem{Schnakenberg1976}
J.~Schnakenberg,
\emph{Network theory of microscopic and macroscopic behavior of master equation systems},
Reviews of Modern Physics, vol.~48, no.~4, pp.~571--585, 1976.
\href{https://doi.org/10.1103/RevModPhys.48.571}{doi:10.1103/RevModPhys.48.571}.

\bibitem{Maas2011}
J.~Maas,
\emph{Gradient flows of the entropy for finite Markov chains},
Journal of Functional Analysis, vol.~261, no.~8, pp.~2250--2292, 2011.
\href{https://doi.org/10.1016/j.jfa.2011.06.009}{doi:10.1016/j.jfa.2011.06.009}.

\bibitem{Noether1971}
E.~Noether,
\emph{Invariant Variation Problems},
Transport Theory and Statistical Physics, vol.~1, no.~3, pp.~186--207, 1971,
English translation of the 1918 original by M.~A. Tavel.
\href{https://doi.org/10.1080/00411457108231446}{doi:10.1080/00411457108231446}.

\bibitem{AharonovBohm1959}
Y.~Aharonov and D.~Bohm,
\emph{Significance of Electromagnetic Potentials in the Quantum Theory},
Physical Review, vol.~115, pp.~485--491, 1959.
\href{https://doi.org/10.1103/PhysRev.115.485}{doi:10.1103/PhysRev.115.485}.

\bibitem{Einstein1916}
A.~Einstein,
\emph{Die Grundlage der allgemeinen Relativit\"atstheorie},
Annalen der Physik, vol.~354, no.~7, pp.~769--822, 1916.
\href{https://doi.org/10.1002/andp.19163540702}{doi:10.1002/andp.19163540702}.

\bibitem{Hawking1975}
S.~W. Hawking,
\emph{Particle creation by black holes},
Communications in Mathematical Physics, vol.~43, pp.~199--220, 1975.
\href{https://doi.org/10.1007/BF02345020}{doi:10.1007/BF02345020}.

\bibitem{GibbonsHawking1977}
G.~W. Gibbons and S.~W. Hawking,
\emph{Action integrals and partition functions in quantum gravity},
Physical Review D, vol.~15, no.~10, pp.~2752--2756, 1977.
\href{https://doi.org/10.1103/PhysRevD.15.2752}{doi:10.1103/PhysRevD.15.2752}.

\bibitem{Bell1964}
J.~S. Bell,
\emph{On the Einstein Podolsky Rosen paradox},
Physics Physique Fizika, vol.~1, no.~3, pp.~195--200, 1964.
\href{https://doi.org/10.1103/PhysicsPhysiqueFizika.1.195}{doi:10.1103/PhysicsPhysiqueFizika.1.195}.

\bibitem{CHSH1969}
J.~F. Clauser, M.~A. Horne, A.~Shimony, and R.~A. Holt,
\emph{Proposed Experiment to Test Local Hidden-Variable Theories},
Physical Review Letters, vol.~23, no.~15, pp.~880--884, 1969.
\href{https://doi.org/10.1103/PhysRevLett.23.880}{doi:10.1103/PhysRevLett.23.880}.

\end{thebibliography}

% ---- END monograph/frontmatter/references.tex ----
